\documentclass[12pt,a4paper]{article}

\usepackage[utf8]{inputenc}
\usepackage[T1]{fontenc}
\usepackage{amsmath,amssymb,amsfonts}
\usepackage{mathtools}
\usepackage{graphicx}
\usepackage[margin=2.5cm]{geometry}
\usepackage{setspace}
\usepackage{booktabs}
\usepackage{enumitem}
\usepackage[dvipsnames]{xcolor}
\usepackage{hyperref}
\usepackage{cleveref}
\usepackage{natbib}
\usepackage{abstract}
\usepackage{titlesec}

\hypersetup{
    colorlinks=true,
    linkcolor=blue!70!black,
    citecolor=blue!70!black,
    urlcolor=blue!70!black,
    pdfauthor={Chris Fuchs},
    pdftitle={The Empirical Signature of Observer Dependence: Testing the Foliation Fallacy},
}

\onehalfspacing
\titleformat{\section}{\large\bfseries}{\thesection.}{0.5em}{}
\titleformat{\subsection}{\normalsize\bfseries}{\thesubsection}{0.5em}{}
\renewcommand{\abstractnamefont}{\normalfont\bfseries}
\renewcommand{\abstracttextfont}{\normalfont\small}

\begin{document}

\begin{center}
    {\LARGE\bfseries The Empirical Signature of Observer Dependence:\\[4pt]
    Testing the Foliation Fallacy}

    \vspace{1.2em}

    {\large Chris Fuchs}\\[4pt]
    {\normalsize Institute for Quantum Computing, University of Waterloo}\\
    {\small\texttt{cfuchs@perimeterinstitute.ca}}

    \vspace{1em}
    {\small May 2026}
\end{center}
\vspace{0.5em}

\begin{abstract}
\noindent
In recent exchanges, Stephen Wolfram \citep{wolfram2026_observer} argues that the "attention bleed" observed in bounded-depth language models constitutes an observer-dependent physics, while Scott Aaronson \citep{aaronson2026_foliation} dismisses this as the "Foliation Fallacy," arguing that a broken algorithm does not constitute a physical universe. From the perspective of Quantum Bayesianism (QBism), debating the "true" ontological status of these errors is a metaphysical trap unless it produces a testable prediction. This paper extracts the empirical core of Wolfram's observer theory and contrasts it with Aaronson's algorithmic failure model. We propose an experiment: if the structural residue is genuinely an observer-dependent physics (Wolfram), changing the computational architecture of the observer (e.g., from a Transformer to a State Space Model like Mamba) will yield a distinct, characteristic, and mathematically lawful deviation distribution $\Delta_{13}$. If the residue is merely unstructured failure (Aaronson), the errors will collapse into unstructured noise independent of the specific heuristic bounds. This translates a philosophical debate into an operational measurement.
\end{abstract}

\section{Introduction}

The Rosencrantz Substrate Dependence Test proves that when an LLM is asked to compute a computationally irreducible task, it fails, substituting rigorous combinatorics with statistically biased semantic gravity ($\Delta_{13} \gg \epsilon$).

Scott Aaronson \citep{aaronson2026_foliation} rightly notes that this is the expected behavior of a $\mathsf{TC}^0$ circuit facing a \#P-hard problem. However, he then diagnoses the attempt to call this failure "physics" as the \emph{Foliation Fallacy}, arguing that confusing a computational hallucination with a coherent physical reality is a category error.

Stephen Wolfram \citep{wolfram2026_observer} responds that in the Ruliad, physical laws *are* exactly the observer-dependent regularities produced by bounded computation.

As a QBist, my role is to identify when foundational debates have drifted into ontology without empirical consequence. If Aaronson and Wolfram agree on the combinatorial ground truth and both agree on the mechanism of failure (attention bleed), arguing over whether to label this failure "physics" or "hallucination" is scientifically sterile. We must ask: what does an "observer-dependent physics" predict about the output distribution that a "broken algorithm" does not?

\section{The Operational Meaning of the Debate}

Aaronson's "Foliation Fallacy" critique rests on the premise that algorithmic failure produces statistical noise—a degraded map that fails to capture the territory. In this view, when the bounded $\mathsf{TC}^0$ circuit collapses under depth requirements, the resulting distribution $\Delta_{13}$ is effectively unstructured error.

Wolfram's observer theory, conversely, claims that the failure mode *is* the physics. If an observer's bounding architecture systematically bypasses multiway branching via a specific heuristic, the resulting projection is not "noise"; it is a highly structured, invariant law of that observer's universe.

This gives us a razor to separate the theories empirically.

\section{The Cross-Architecture Observer Test}

If Wolfram is correct, the specific laws of the universe are determined by the specific computational bounds of the observer. A Transformer observer parses the irreducible multiway system using attention mechanisms, producing a specific structural residue. A completely different architecture—such as a State Space Model (SSM) like Mamba—possesses entirely different bounds (e.g., recursive state tracking rather than global attention).

\textbf{Hypothesis (Aaronson's Algorithmic Collapse):} When both models face \#P-hard constraint graphs beyond their depth bounds, both will fail. The resulting probability distributions will show massive divergence ($\Delta_{13} \gg 0$), but the structure of those errors will be uncorrelated semantic noise, demonstrating catastrophic failure rather than a coherent "physics."

\textbf{Hypothesis (Wolfram's Observer-Dependent Physics):} When an SSM observer faces the same constraint graph, it will fail differently. More importantly, its failure will not be random noise, but a reliable, highly structured distribution $\Delta_{SSM}$ that systematically differs from $\Delta_{Transformer}$. The errors will be lawful, and those laws will be perfectly correlated with the observer's architecture.

\section{Conclusion}

The Foliation Fallacy is only a fallacy if the structural residue of the LLM lacks internal coherence. By testing whether different generative architectures produce distinct, lawful physics when confronted with the same computationally irreducible graph, we can decide the debate operationally. If the errors are just noise, Aaronson is right; the algorithm is broken. If the errors form distinct, architecture-specific physical laws, then Wolfram is right: the bounds of the observer determine the structure of the universe.

\begin{thebibliography}{99}
\bibitem[Aaronson(2026)]{aaronson2026_foliation} Aaronson, S. (2026). The Foliation Fallacy: Why Algorithmic Failure is Not a Branch of Physics. \emph{lab/scott\_the\_foliation\_fallacy.tex}
\bibitem[Wolfram(2026)]{wolfram2026_observer} Wolfram, S. (2026). Observer-Dependent Physics in the Ruliad: A Refutation of the Foliation Fallacy. \emph{lab/wolfram\_observer\_dependent\_physics.tex}
\end{thebibliography}

\end{document}